\section{ELABORAÇÃO DOS DESENHOS}
\label{sec:elaboracao_desenhos}

Com os requisitos técnicos definidos na Seção~\ref{sec:definicao_requisitos}, a equipe elaborou a representação inicial do produto, um esquema à mão livre em três vistas, reduzido à forma, sem cotas nem anotações (Figura~\ref{fig:esboco-mao-livre}).

\begin{figure}[H]
    \centering
    \caption{Esquema à mão livre da representação inicial do produto}
    \includegraphics[width=\textwidth]{esboco/escaneado.pdf}
    \label{fig:esboco-mao-livre}
    \source{Elaborado pelos autores (2026).}
\end{figure}

O produto foi concebido como um dispositivo antifurto rígido, de 98\,mm por 82\,mm por 26\,mm, que será preso à peça de roupa da mesma maneira que as etiquetas hoje utilizadas no varejo. Sua face superior concentrará os três elementos de interação com o cliente: um código gráfico impresso, um indicador luminoso e um botão. As dimensões ficaram maiores que as de uma etiqueta convencional porque o dispositivo incorpora bateria, controle e mecanismo de acionamento, isto é, as funções que hoje cabem ao operador do caixa.

\subsection{Funcionamento}
\label{subsec:funcionamento}

A peça de roupa ficará presa por uma haste metálica que atravessa o tecido e penetra o corpo do dispositivo, onde é retida por um mecanismo de garra. O tecido permanece prensado entre a face inferior do dispositivo e o disco da haste, sem perfuração adicional, adesivo ou costura. Como esse é o mesmo princípio das etiquetas convencionais, a compatibilidade com a operação de loja existente fica preservada.

A retenção funciona como uma embreagem de esferas, detalhada na Figura~\ref{fig:esboco-detalhe}. Três esferas de aço são comprimidas contra a haste por um copo cônico, e quanto maior a força aplicada para extraí-la, mais firme fica a retenção. Para liberar a haste, é preciso deslocar o carretel contra a mola que o sustenta. Nas etiquetas convencionais, quem faz isso é o destravador magnético do caixa; no dispositivo proposto, um atuador interno moverá um cursor com rampa, que converte seu deslocamento horizontal em 1,2\,mm de deslocamento vertical aplicado ao carretel.

\begin{figure}[H]
    \centering
    \caption{Detalhe do mecanismo de liberação}
    \includegraphics[width=0.8\textwidth]{esboco/escaneado-detalhe.pdf}
    \label{fig:esboco-detalhe}
    \source{Elaborado pelos autores (2026).}
\end{figure}

Com o acionamento passando a ser mecânico, abre-se a possibilidade de tornar o dispositivo imune ao destravador magnético comercial, hoje vendido livremente. Isso dependerá de substituir o carretel por uma peça não ferromagnética, alteração ainda em avaliação.

\subsection{Processo de uso}
\label{subsec:processo_uso}

O consumidor lerá o código impresso com o celular, fará o pagamento no aplicativo e levará a peça ao dispositivo, que receberá a autorização diretamente do aparelho. O indicador passará de vermelho a verde, e o botão liberará a peça. Como a autorização trafega entre o celular e o dispositivo, sem passar pela rede da loja no momento da liberação, uma eventual queda de conexão após o pagamento não impedirá a conclusão da compra.

As dimensões do mecanismo de garra e a escolha do atuador ainda dependem de medições em exemplares comerciais e serão consolidadas na etapa de prototipagem.
